% Procedencia: integral 2249, cap. 24, pp. 609--618; capitulo_24_dualidad_t.tex.
% Las fórmulas de dominio operatorio se explicitan sobre la realización ya construida.
\section{Dualité T : coordonnées entières, changement de section et équivalence opératorielle}
\label{iv:dualidad}

La décomposition arithmétique de APP conserve le représentant résiduel
et son quotient. L’évolution TPK transporte en outre son orientation et
la mémoire des franchissements de bord. L’échange des deux
coordonnées exige de conserver cette distinction et de préciser le domaine
dans lequel l’opération est de nouveau définie. Sur ce domaine, on
construit d’abord une involution arithmétique ; sa forme quadratique
admet ensuite une représentation par des opérateurs autoadjoints.

\subsection{Domaine stable de l’échange}

Dans la section résiduelle positive,
\begin{equation}
 r_9^+(n)=1+((n-1)\bmod9),\qquad
 Q_9^+(n)=\frac{n-r_9^+(n)}9,
 \qquad n=r_9^+(n)+9Q_9^+(n).
 \label{iv:eq:division-positiva}
\end{equation}
Un échange peut placer un quotient arbitraire dans la première
coordonnée. Son domaine stable est donc
\(\mathscr D_0=\ZZ^2\times\RR_{>0}\), et non seulement la section
\(1\leq r\leq9\). Le prolongement du domaine permet de définir
une opération totale ; l’admissibilité d’une paire comme lecture d’une
route TPK conserve séparément ses règles de sélection.

\begin{definition}[Forme compacte normalisée et échange
]
\label{iv:def:dualidad}
Pour \(d=(m,w,R_0)\in\mathscr D_0\), soient
\begin{equation}
 E_{R_0}(m,w)=\frac{m^2}{R_0^2}+w^2R_0^2,
 \qquad
 \mathsf T_0(m,w,R_0)=(w,m,R_0^{-1}).
 \label{iv:eq:energia-t}
\end{equation}
Les symboles \(m,w\) nomment les deux coordonnées entières
transportées ; \(R_0\) est l’échelle positive sans dimension de
la forme. Leur interprétation au moyen de l’impulsion, de l’enroulement et du
rayon physique est établie par une application de réalisation ultérieure.
\end{definition}

\begin{theorem}[Involution et conservation de la forme quadratique
]
\label{iv:thm:dualidad-escalar}
L’application \(\mathsf T_0\) est une involution de
\(\mathscr D_0\), et

\[
 E_{R_0^{-1}}(w,m)=E_{R_0}(m,w).
\]
\end{theorem}
\begin{proof}
Deux échanges restituent \((m,w)\), et deux inversions restituent
\(R_0\). De plus,
\[
 E_{R_0^{-1}}(w,m)=w^2R_0^2+\frac{m^2}{R_0^2}
 =E_{R_0}(m,w).
\]
\end{proof}

L’équivalence vaut pour chaque rayon et son réciproque. En \(R_0=1\),
le rayon reste fixe et l’opération agit sur une même forme.
Une paire concrète est fixe par l’involution complète précisément
lorsque, en outre, \(m=w\).


\subsection{Hamiltonien compact et domaine auto-adjoint}

Soit \(\mathcal H_{\mathrm{lat}}=\ell^2(\ZZ^2)\), de base
orthonormale \(\{|m,w\rangle\}\). Nous définissons
\begin{align}
 H(R_0)|m,w\rangle&=E_{R_0}(m,w)|m,w\rangle,\label{iv:eq:hamiltoniano}\\
 D(H(R_0))&=\left\{\psi:\sum_{m,w\in\ZZ}
 E_{R_0}(m,w)^2|\psi_{m,w}|^2<\infty\right\}.\label{iv:eq:dominio-h}
\end{align}
Les vecteurs à support fini appartiennent au domaine et forment un
cœur de l’opérateur diagonal. La forme quadratique associée a pour
domaine
\[
 D(H(R_0)^{1/2})=\left\{\psi:\sum_{m,w}
 E_{R_0}(m,w)|\psi_{m,w}|^2<\infty\right\}.
\]

\begin{proposition}[Caractère autoadjoint et résolvante compacte]
\label{iv:prop:h-autoadjunto}
L’opérateur \(H(R_0)\) est autoadjoint et non négatif. Ses
valeurs propres sont \(E_{R_0}(m,w)\), comptées avec multiplicité,
et \((H(R_0)+I)^{-1}\) est compact.
\end{proposition}
\begin{proof}
Un opérateur de multiplication par une suite réelle est symétrique
sur \eqref{iv:eq:dominio-h}. Si \(\eta\) appartient au domaine
de l’adjoint, tester l’identité de l’adjoint contre chaque vecteur
\(|m,w\rangle\) impose que son image ait pour coordonnées
\(E_{R_0}(m,w)\eta_{m,w}\). La condition d’appartenance à
\(\ell^2\) est exactement \eqref{iv:eq:dominio-h} ; l’adjoint a donc
le même domaine et la même action. La non-négativité
s’obtient en additionnant les poids non négatifs.

Avec \(c_{R_0}=\min(R_0^{-2},R_0^2)>0\), on a
\(E_{R_0}(m,w)\geq c_{R_0}(m^2+w^2)\). Hors de tout
ensemble fini croissant, les coefficients
\((1+E_{R_0}(m,w))^{-1}\) tendent uniformément vers zéro.
Tronquer la résolvante sur ces ensembles produit des opérateurs
de rang fini qui convergent en norme, ce qui prouve sa
compacité. La base diagonale détermine le spectre annoncé.

\end{proof}

\begin{theorem}[Conjugaison unitaire, domaines compris]
\label{iv:thm:t-unitaria}
\label{prop:dualidad-radios-k}
L’échange
\(U_T|m,w\rangle=|w,m\rangle\) est unitaire,
\(U_T^2=I\), et satisfait
\begin{equation}
 U_TD(H(R_0))=D(H(R_0^{-1})),
 \qquad U_TH(R_0)U_T^{-1}=H(R_0^{-1}).
 \label{iv:eq:conjugacion-h}
\end{equation}
Il transporte aussi les domaines des formes quadratiques.
\end{theorem}
\begin{proof}
L’échange est une permutation involutive de la base
orthonormale. Pour une suite \(\psi\), l’identité du
théorème \ref{iv:thm:dualidad-escalar} donne
\[
 \sum_{m,w}E_{R_0^{-1}}(m,w)^2|(U_T\psi)_{m,w}|^2
 =\sum_{m,w}E_{R_0}(m,w)^2|\psi_{m,w}|^2.
\]
On obtient ainsi l’égalité des domaines. Sur chaque coordonnée,
\(U_TH(R_0)U_T^{-1}\) multiplie par
\(E_{R_0}(w,m)=E_{R_0^{-1}}(m,w)\). La même somme, avec
la première puissance du poids, démontre le résultat pour les
formes quadratiques.
\end{proof}

L’égalité scalaire et l’égalité opératorielle sont deux réalisations
de la même involution. La seconde conserve les domaines qui
permettent d’interpréter la première comme une équivalence spectrale,
non seulement comme une égalité sur des vecteurs à support fini.

\subsection{Spécialisation au rayon de l’ellipse d’action}

Le rayon \(R_\star\) déterminé dans la section
\ref{iv:ellipse:section} fixe la réalisation de cette collection qui
correspond aux représentations angulaires de l’état. Le symbole
\(R_0\) permet d’exprimer les résultats de transport pour toute la
collection positive ; la composition avec l’ellipse évalue cette collection en
\(R_0=R_\star\) et conserve aussi son image réciproque.

\begin{corollary}[Collection elliptique stable par dualité]
\label{iv:cor:radio-especializado}
L'ensemble
\[
 \mathscr D_{\mathrm{el}}
 =\ZZ^2\times\{R_\star,R_\star^{-1}\}
\]
est stable par \(\mathsf T_0\). Ses hamiltoniens sont autoadjoints,
positifs ou nuls et à résolvante compacte. L’échange \(U_T\)
transporte le domaine de \(H(R_\star)\) sur celui de
\(H(R_\star^{-1})\), y compris les domaines de leurs formes.

\end{corollary}
\begin{proof}
La positivité de \(R_\star\) et sa détermination unique sont établies
dans la section \ref{iv:ellipse:section}. L’inversion permute
\(R_\star\) et \(R_\star^{-1}\), tandis que l’échange préserve
\(\ZZ^2\). Dans ces deux fibres, on applique la proposition
\ref{iv:prop:h-autoadjunto} et le théorème \ref{iv:thm:t-unitaria}.
Si \(R_\star=1\), les deux fibres coïncident ; dans tout autre cas,
l’opérateur entrelace deux fibres distinctes de la même collection.
\end{proof}

L’échelle commune d’action détermine l’aire de l’ellipse ; le quotient
de ses demi-axes détermine \(R_\star\). Cette séparation permet de conserver
l’aire et l’échelle commune d’action en échangeant les demi-axes et d’obtenir la réciprocité
du rayon. L’échange d’indices sur \(\ell^2(\ZZ^2)\)
et les transports du plan d’action agissent dans des espaces différents ;
leurs propriétés respectives sont démontrées dans leurs domaines propres.


\subsection{Changement de représentant et mémoire du quotient}

La section zéro utilise \(\bar r\in\{0,\ldots,8\}\). Le changement
qui conserve l’entier de \eqref{iv:eq:division-positiva} est
\begin{equation}
 C(r,Q)=(\bar r,Q+\mathbf1_{\{r=9\}}).
 \label{iv:eq:cambio-seccion}
\end{equation}
La modification du quotient fait partie de l’opération.
Pour \(n=9\), les deux paires sont \((9,0)\) et \((0,1)\).
Leurs énergies non corrigées sont \(81/R_0^2\) et \(R_0^2\).
Représenter le même entier n’implique donc pas l’égalité de
ces deux insertions dans une forme qui distingue les coordonnées.

La relation d’équivalence est conservée en définissant
\[
 L_9=\ZZ(9,-1),\qquad S(x_1,x_2)=(x_2,x_1),
 \qquad q_{R_0}(x)=\frac{x_1^2}{R_0^2}+x_2^2R_0^2.
\]
Pour \(x\in\ZZ^2\), soit
\begin{equation}
 \mathcal E_{R_0}^{L_9}([x])
 =\min_{k\in\ZZ}q_{R_0}(x+k(9,-1)).
 \label{iv:eq:energia-cociente}
\end{equation}
Cette construction opère sur des classes. La mémoire qui distingue
les représentants individuels reste dans l’état TPK antérieur
au quotient ; la fonction de \eqref{iv:eq:energia-cociente}
détermine l’évaluation invariante de la classe choisie.

\begin{theorem}[Dualité covariante des sections
]
\label{iv:thm:dualidad-cociente}
L’énergie de \eqref{iv:eq:energia-cociente} est bien définie.
La transformation

\[
 \mathsf T_{\mathrm{cov}}([x]_{L_9},R_0)
 =([Sx]_{SL_9},R_0^{-1})
\]
est une bijection entre
\((\ZZ^2/L_9)\times\RR_{>0}\) et
\((\ZZ^2/SL_9)\times\RR_{>0}\), dont l’inverse est donné par le
second échange, et
\[
 \mathcal E_{R_0^{-1}}^{SL_9}([Sx])
 =\mathcal E_{R_0}^{L_9}([x]).
\]
\end{theorem}
\begin{proof}
Remplacer \(x\) par \(x+k_0(9,-1)\) réindexe les candidats
au minimum. Celui-ci existe : la fonction de \(k\) est quadratique et
a pour coefficient dominant \(81/R_0^2+R_0^2>0\). Le changement
\((9,Q)\mapsto(0,Q+1)\) a une différence dans \(L_9\), de
sorte qu’il représente une même classe. L’échange transporte
la relation d’équivalence vers \(SL_9\), et
\[
 \min_{\ell\in L_9}q_{R_0^{-1}}(Sx+S\ell)
 =\min_{\ell\in L_9}q_{R_0}(x+\ell).
\]
Enfin, \(S^2=I\) récupère le quotient et le rayon initiaux.

\end{proof}

L’évaluation du minimum peut être effectuée exactement avec deux
candidats. Si \(x=(m,w)\), le minimum réel de la quadratique
se produit en
\[
 k_* =\frac{wR_0^4-9m}{81+R_0^4}.
\]
Le minimum entier est atteint en \(\lfloor k_*\rfloor\) ou en
\(\lceil k_*\rceil\). Cette formule n’altère pas l’équivalence ;
elle explicite son évaluation finie et permet de vérifier l’effet
du changement de section sans écarter le quotient.


\subsection{Normalisation du rayon dans la réalisation de corde}

Soit \(\alpha'>0\) une échelle de la réalisation de corde.
Ce symbole désigne l’échelle de corde et reste distinct
de la constante de structure fine \(\alpha\). Nous définissons
\begin{equation}
 \mathsf T_{\alpha'}(m,w,R)
 =\left(w,m,\frac{\alpha'}R\right),\qquad
 \mathcal N_{\alpha'}(m,w,R)
 =\left(m,w,\frac R{\sqrt{\alpha'}}\right).
 \label{iv:eq:normalizacion-cuerda}
\end{equation}

\begin{theorem}[Conjugaison exacte par l’échelle]
\label{iv:thm:normalizacion-cuerda}
La normalisation satisfait

\[
 \mathcal N_{\alpha'}\mathsf T_{\alpha'}
 =\mathsf T_0\mathcal N_{\alpha'}.
\]
Si
\[
 E_{\mathrm{str}}(m,w,R)
 =\frac{m^2}{R^2}+\frac{w^2R^2}{\alpha'^2},
\]
alors
\(E_{\mathrm{str}}=\alpha'^{-1}E\circ\mathcal N_{\alpha'}\).
Les combinaisons

\[
 p_L=\frac mR+\frac{wR}{\alpha'},\qquad
 p_R=\frac mR-\frac{wR}{\alpha'}
\]
se transforment au moyen de \(p_L\mapsto p_L\) et \(p_R\mapsto-p_R\).
\end{theorem}
\begin{proof}
Les deux compositions de la première égalité donnent
\((w,m,\sqrt{\alpha'}/R)\). Par ailleurs,

\[
 E_{R/\sqrt{\alpha'}}(m,w)
 =\frac{\alpha'm^2}{R^2}+\frac{w^2R^2}{\alpha'}
 =\alpha'E_{\mathrm{str}}(m,w,R).
\]
Substituer \((w,m,\alpha'/R)\) dans les expressions de
\(p_L,p_R\) donne, respectivement, les deux transformations
affirmées.
\end{proof}

La conjugaison fixe ce que signifie exprimer le même résultat dans
un système de coordonnées dimensionnelles. Dans la collection elliptique, les rayons sont
\[
 R=\sqrt{\alpha'}R_\star,\qquad
 R^\vee=\sqrt{\alpha'}R_\star^{-1}.
\]
Leur produit est \(\alpha'\). La détermination interne du quotient
adimensionnel des demi-axes et l’échelle de longueur au carré de la
réalisation remplissent des fonctions distinctes : la première fixe la forme ;
la seconde exprime ses rayons dans une unité dimensionnelle.

La normalisation précédente permet d’exprimer le même résultat en
rayon normalisé ou physique. L’attribution d’une tension et d’une
unité à \(\alpha'\) appartient à l’application de réalisation et a
sa propre origine ; ce n’est pas une donnée intervenant dans la production
des registres APP--TRIT--TPK.

\subsection{Relation avec la transformation modulaire
}

L’inversion du rayon a une représentation précise sur
l’axe imaginaire du demi-plan supérieur. Soient
\(\iota(R_0)=\mathrm iR_0^2\) et
\(\mathsf S(\tau)=-1/\tau\). Alors

\begin{equation}
 \mathsf S\circ\iota(R_0)=\iota(R_0^{-1}).
 \label{iv:eq:semiconjugacion-modular}
\end{equation}
L’égalité résulte de
\(-1/(\mathrm iR_0^2)=\mathrm iR_0^{-2}\). Pour la fonction
modulaire \(J=j-744\), dont l’invariance est conservée dans le
développement de Moonshine, on déduit
\(J(\mathrm iR_0^2)=J(\mathrm iR_0^{-2})\).

L’application \(\iota\) exprime une semi-conjugaison de paramètres.
L’échange \(U_T\) agit sur les coordonnées de
\(\ell^2(\ZZ^2)\) ; \(\mathsf S\) agit sur le demi-plan
supérieur ; l’automorphisme dual d’un orbifold agit dans ses
secteurs gradués. L’écriture de leurs domaines permet de
composer des relations spécifiques entre eux en conservant
l’information de chaque représentation. Les théorèmes
\ref{iv:thm:t-unitaria}, \ref{iv:thm:dualidad-cociente}
et \ref{iv:thm:normalizacion-cuerda} établissent, respectivement,
la conjugaison unitaire, le changement de section résiduelle et la
normalisation d’échelle, avec les domaines opératoriels déclarés.
